% Black–Scholes & the Greeks — research-note PDF.
%
% PROSE AND CODE ARE VERBATIM from the website article:
%   site/src/content/articles/black-scholes-greeks.tsx
%   site/src/content/articles/black-scholes-greeks-code.ts
% Metadata from lib/articles.ts; format and project-card fallbacks from
% lib/topics.ts (the registry entry sets neither).
%
% CHARTS generated by figures/make_black_scholes_greeks.py from the article's
% own data module (data/black-scholes-and-the-greeks.ts) — the same computed
% results the site's <Line> components render. Figure captions are the
% article's own, which already carry its numbering ("Figure 4.2 · ...").
\documentclass{pyportfolios-note}

\noteshorttitle{Black–Scholes \& the Greeks}

\begin{document}
\thispagestyle{notecover}

\notemasthead{Quantitative Insights}{Tutorial}{Q3 2026}{}

\notetitle{Black–Scholes\\\& the Greeks}

\notedeck{The BSM model turns five observable inputs into a fair option price and
a full risk report — shown through QQQ (Nasdaq 100) options.}

\notelede{The Black–Scholes–Merton (BSM) model turns five observable inputs —
spot, strike, time, rate, volatility — into a fair option price \textit{and} a
full risk report (the Greeks). Fifty years on, it remains the language in which
options are quoted, hedged, and risk-managed.}

\notecard
  {Quant Foundations \& Derivatives}
  {NumPy · SciPy · Matplotlib}
  {QQQ options (Nasdaq 100 ETF, Invesco)}
  {Jan 2018 – Dec 2024 (trailing 1y vol)}
  {Pricing \& hedging options; implied vol is how the entire options market
   quotes prices}

% ============================================================= 1 Summary =====
\noteneedroom{150bp}
\notesection{1.\notenumsep Summary}

\begin{multicols}{2}
BSM assumes the underlying follows \textbf{Geometric Brownian Motion} (see the
previous tutorial) and shows that an option's payoff can be \textbf{replicated}
by continuously trading the stock and a bond — so the option's price is the cost
of that replication, independent of anyone's market view. The result is a
closed-form price for European calls and puts, and analytic \textbf{Greeks}: the
sensitivities of that price to spot ($\Delta$, $\Gamma$), volatility ($\nu$,
``vega''), time ($\Theta$) and rates ($\rho$).
\end{multicols}

% =========================================================== 2 Intuition =====
\noteneedroom{170bp}
\notesection{2.\notenumsep Intuition}

\begin{multicols}{2}
\notesubsection{2.1\notenumsep Replication, not prediction}
Hold $\Delta$ shares against a short call and the portfolio is (momentarily)
immune to small price moves. Rebalance continuously and the option is
manufactured from stock and cash. No-arbitrage then pins the price — the
expected return of the stock $\mu$ \textbf{drops out entirely}. That is the
magic: two people who disagree wildly about where the Nasdaq is going must still
agree on the option's fair price.

\notesubsection{2.2\notenumsep Volatility is the price of uncertainty}
Of the five inputs, four are observable. Volatility is not — and the option
premium is, in essence, the market's bid on future uncertainty. Higher $\sigma$
$\rightarrow$ wider terminal distribution $\rightarrow$ both calls and puts gain
value (optionality only benefits from dispersion).

\notesubsection{2.3\notenumsep Greeks: the risk dashboard}
A trader rarely asks ``what is the price?'' — they ask ``what happens to my book
if spot drops 1\%, vol jumps 2 points, and a day passes?'' Delta, Gamma, Vega
and Theta answer exactly that, one partial derivative at a time.
\end{multicols}

\begin{notetable}{@{}p{120bp}p{170bp}p{206bp}@{}}
\thcell{Greek} & \thcell{Sensitivity to} & \thcell{Typical use} \\
\theadrule
Delta $\Delta$ & Spot & Hedge ratio \\ \rowrule
Gamma $\Gamma$ & Spot (2nd order) & Hedge stability / convexity \\ \rowrule
Vega $\nu$ & Volatility & Vol exposure \\ \rowrule
Theta $\Theta$ & Time & Daily carry / decay \\ \rowrule
Rho $\rho$ & Interest rate & Rate exposure \\
\end{notetable}

% ================================================= 3 Theory & Mechanics ======
\noteneedroom{190bp}
\notesection{3.\notenumsep Theory \& Mechanics}

\notesubsection{3.1\notenumsep The pricing formula}

With spot $S$, strike $K$, maturity $T$, rate $r$ and volatility $\sigma$:

\[ d_1 = \frac{\ln(S/K) + (r + \tfrac{1}{2}\sigma^2)T}{\sigma\sqrt{T}},
   \qquad d_2 = d_1 - \sigma\sqrt{T} \]

\[ C = S\,N(d_1) - K e^{-rT} N(d_2), \qquad
   P = K e^{-rT} N(-d_2) - S\,N(-d_1) \]

Read $N(d_2)$ as the (risk-neutral) probability the option finishes in the
money; $N(d_1)$ is the delta of the call.

\notesubsection{3.2\notenumsep The Greeks in closed form}

\[ \Delta_C = N(d_1), \quad \Gamma = \frac{N'(d_1)}{S\sigma\sqrt{T}}, \quad
   \nu = S\,N'(d_1)\sqrt{T}, \quad
   \Theta_C = -\frac{S N'(d_1)\sigma}{2\sqrt{T}} - rKe^{-rT}N(d_2) \]

\begin{notecode}
def bs_price(S, K, T, r, sigma, kind="call"):
    """Black-Scholes price for a European call or put."""
    d1 = (np.log(S / K) + (r + 0.5 * sigma**2) * T) / (sigma * np.sqrt(T))
    d2 = d1 - sigma * np.sqrt(T)
    if kind == "call":
        return S * norm.cdf(d1) - K * np.exp(-r * T) * norm.cdf(d2)
    return K * np.exp(-r * T) * norm.cdf(-d2) - S * norm.cdf(-d1)

def bs_greeks(S, K, T, r, sigma, kind="call"):
    """Delta, Gamma, Vega, Theta (per year), Rho for a European option."""
    d1 = (np.log(S / K) + (r + 0.5 * sigma**2) * T) / (sigma * np.sqrt(T))
    d2 = d1 - sigma * np.sqrt(T)
    sign = 1 if kind == "call" else -1
    return {
        "delta": sign * norm.cdf(sign * d1),
        "gamma": norm.pdf(d1) / (S * sigma * np.sqrt(T)),
        "vega":  S * norm.pdf(d1) * np.sqrt(T),
        "theta": (-S * norm.pdf(d1) * sigma / (2 * np.sqrt(T))
                  - sign * r * K * np.exp(-r * T) * norm.cdf(sign * d2)),
        "rho":   sign * K * T * np.exp(-r * T) * norm.cdf(sign * d2),
    }
\end{notecode}

\notesubsection{3.3\notenumsep Put-call parity — the free sanity check}

Independent of any model: $C - P = S - Ke^{-rT}$. If your implementation
violates it, something is wrong.

% ============================================ 4 Applied Example — QQQ ========
\noteneedroom{160bp}
\notesection{4.\notenumsep Applied Example — QQQ}

\notesubsection{4.1\notenumsep Inputs from real data}

Spot from the latest QQQ close; volatility proxied by the trailing 1-year
realised vol of log returns (in practice you would use the implied vol quoted in
the market — that distinction is exactly what the \textit{Heston vs
Black-Scholes} piece explores).

\begin{notecode}
TICKER = "QQQ"
START, END = "2018-01-01", "2024-12-31"

def load_prices(ticker, start, end):
    """Adjusted-close prices: yfinance first, Stooq as fallback."""
    try:
        import yfinance as yf
        df = yf.download(ticker, start=start, end=end, auto_adjust=True, progress=False)
        if not df.empty:
            return df["Close"].squeeze().rename(ticker).dropna()
    except Exception as exc:
        print(f"yfinance failed ({exc}); trying Stooq…")
    url = f"https://stooq.com/q/d/l/?s={ticker.lower()}.us&i=d"
    df = pd.read_csv(url, parse_dates=["Date"], index_col="Date")
    return df.loc[start:end, "Close"].rename(ticker).dropna()

px = load_prices(TICKER, START, END)
log_ret = np.log(px / px.shift(1)).dropna()

S0    = float(px.iloc[-1])
SIGMA = float(log_ret.tail(252).std() * np.sqrt(252))   # trailing 1y realised vol
R     = 0.045                                           # short-term risk-free rate
T3M   = 0.25                                            # 3 months
K_ATM = round(S0)                                       # at-the-money strike

call = bs_price(S0, K_ATM, T3M, R, SIGMA, "call")
put  = bs_price(S0, K_ATM, T3M, R, SIGMA, "put")

print(f"{TICKER} spot = {S0:,.2f}   sigma = {SIGMA:.2%}   r = {R:.2%}")
print(f"3M ATM call (K={K_ATM}): {call:6.2f}")
print(f"3M ATM put  (K={K_ATM}): {put:6.2f}")

g = bs_greeks(S0, K_ATM, T3M, R, SIGMA, "call")
print("\n3M ATM call Greeks:")
print(f"  delta {g['delta']:7.3f}   gamma {g['gamma']:8.5f}   vega {g['vega']:7.2f}"
      f"   theta {g['theta']/365:7.3f}/day   rho {g['rho']:6.2f}")
\end{notecode}

\notesubsection{4.2\notenumsep The Greeks across strikes}

The same option book looks completely different in-, at- and out-of-the-money.
Gamma and Vega \textbf{peak at the money} — that is where hedges churn fastest
and vol exposure is largest; Delta rolls from 0 to 1 like a smoothed step
function.

\begin{notecode}
strikes = np.linspace(0.80 * S0, 1.20 * S0, 200)
G = {k: [] for k in ["delta", "gamma", "vega", "theta"]}
for K in strikes:
    gk = bs_greeks(S0, K, T3M, R, SIGMA, "call")
    for k in G: G[k].append(gk[k])

fig, axes = plt.subplots(2, 2, figsize=(11, 6), sharex=True)
for ax, k in zip(axes.flat, G):
    ax.plot(strikes / S0, G[k], color="steelblue")
    ax.axvline(1.0, color="gray", ls=":", lw=1)
    ax.set_title(k.capitalize()); ax.set_xlabel("Moneyness  K / S")
fig.suptitle(f"{TICKER} 3M call Greeks across strikes  (S={S0:,.0f}, sigma={SIGMA:.0%})")
plt.tight_layout(); plt.show()
\end{notecode}

\notefigure{figures/black-scholes-greeks/fig1-greeks.png}
  {Figure 4.2 · QQQ 3-month call — the Greeks across strikes}
  {Delta falls smoothly from 1 to 0 while Gamma, scaled by 100 to share the
   axis, peaks near the money}
  {Source: Author calculations, from the article's computed results. Three-month
   QQQ call, $S$ = 507.45, $\sigma$ = 18.04\%, $r$ = 4\%.}

\notesubsection{4.3\notenumsep Time decay — the option's ticking clock}

Repricing the ATM call as maturity shrinks shows Theta at work: decay is slow far
from expiry and \textbf{accelerates sharply in the final weeks} — the reason
short-dated option selling and 0DTE trading are games of Theta and Gamma.

\begin{notecode}
maturities = np.linspace(1e-3, 1.0, 200)
atm_prices  = [bs_price(S0, K_ATM, t, R, SIGMA, "call") for t in maturities]
otm_prices  = [bs_price(S0, 1.05 * S0, t, R, SIGMA, "call") for t in maturities]

fig, ax = plt.subplots(figsize=(10, 4))
ax.plot(maturities * 12, atm_prices, label=f"ATM (K={K_ATM})", color="steelblue")
ax.plot(maturities * 12, otm_prices, label="5% OTM", color="darkorange")
ax.set_xlabel("Months to expiry"); ax.set_ylabel("Call value")
ax.set_title(f"{TICKER}: value vs time to expiry — Theta accelerates near zero")
ax.invert_xaxis(); ax.legend()
plt.tight_layout(); plt.show()
\end{notecode}

\notefigure{figures/black-scholes-greeks/fig2-decay.png}
  {Figure 4.3 · QQQ call value vs time to expiry — Theta accelerates near zero}
  {At-the-money and 5\% out-of-the-money calls; both decay toward zero and
   steepen in the final weeks}
  {Source: Author calculations, from the article's computed results. Same
   parameters as above.}

\notesubsection{4.4\notenumsep Sanity checks: parity and Monte Carlo}

Two independent tests of the implementation:

\begin{notebullets}
\item \textbf{Put-call parity} must hold to machine precision.
\item \textbf{Monte Carlo under GBM} (the simulator from the previous tutorial,
      with drift $r$) must converge to the closed-form price — Black-Scholes
      \textit{is} the GBM expectation in disguise.
\end{notebullets}

\begin{notecode}
# 1) put-call parity
lhs, rhs = call - put, S0 - K_ATM * np.exp(-R * T3M)
print(f"parity: C - P = {lhs:.6f}   S - K e^-rT = {rhs:.6f}   diff = {abs(lhs-rhs):.2e}")

# 2) Monte Carlo pricing under risk-neutral GBM
N = 200_000
z = np.random.normal(size=N)
S_T = S0 * np.exp((R - 0.5 * SIGMA**2) * T3M + SIGMA * np.sqrt(T3M) * z)
mc_call = np.exp(-R * T3M) * np.maximum(S_T - K_ATM, 0).mean()
se      = np.exp(-R * T3M) * np.maximum(S_T - K_ATM, 0).std() / np.sqrt(N)
print(f"Monte Carlo call: {mc_call:.3f} ± {2*se:.3f}   closed form: {call:.3f}")
\end{notecode}

% ========================================================== 5 Conclusion =====
\noteneedroom{330bp}
\notesection{5.\notenumsep Conclusion}

\begin{multicols}{2}
\notesubsection{5a.\notenumsep Strengths}
\begin{notebullets}
\item \textbf{Closed form} — instant prices and Greeks, no simulation needed
\item \textbf{Preference-free} — the stock's expected return drops out; only
      volatility matters
\item \textbf{The market's lingua franca} — options are \textit{quoted} in BSM
      implied vol even where the model's assumptions fail
\item \textbf{Analytic Greeks} — the entire hedging industry runs on them
\end{notebullets}

\notesubsection{5b.\notenumsep Weaknesses \& Limitations}
\begin{notebullets}
\item \textbf{Constant volatility} — markets trade a \textit{smile/skew}, not a
      flat vol (one number cannot price all strikes)
\item \textbf{GBM underlying} — no jumps, no fat tails; deep OTM puts are
      systematically underpriced
\item \textbf{European exercise only} — American-style early exercise needs
      trees or numerical methods
\item \textbf{Continuous frictionless hedging} — real rebalancing is discrete
      and costs money
\end{notebullets}

\notesubsection{5c.\notenumsep Applications in Practice}
\begin{notebullets}
\item Pricing and hedging European options and warrants
\item Implied volatility extraction — the quoting convention for the entire
      options market
\item Real-time Greek dashboards for market-making and portfolio risk
\end{notebullets}

\notesubsection{5d.\notenumsep Alternatives \& Extensions}
\begin{notebullets}
\item \textbf{Stochastic volatility} — Heston (see our Quant Insights piece
      \textit{``Heston vs Black-Scholes: fitting the volatility smile''})
\item \textbf{Jump-diffusion} — Merton, for crash risk and short-dated skew
\item \textbf{Binomial trees} — American exercise, dividends, path-dependence
\item \textbf{Local volatility} — Dupire, fitting the entire observed surface
\end{notebullets}
\end{multicols}

\end{document}
